\section{A perturbed center with very few terms}\label{cl:sec:center-lift}
\begin{theorem}\label{cl:thm:main}
For each pair of positive integers \(s,d\) there is an integer \(J_0(s,d)\) such that, for every integer \(J\ge J_0(s,d)\), the field \eqref{cl:eq:field} at scale \(t=1/J\) has exactly \(2s+d+1\) nonzero terms and at least \(4sd\) simultaneous hyperbolic cycles. Its cleared degree is \(B=1+L\max(4,2s+1,2d+1)\), where \(L=\Mclr\) is the even clearing factor defined below.
\end{theorem}
The maximum comes from the untilted extremal powers. For \(d=1\) the term giving \(2d+1\) is absent and the maximum is unchanged. Optimizing \(4sd\) over \(2s+d+1=N\) gives \(4,8,12,16,24,32,40\) cycles for \(N=4,\ldots,10\), respectively. These realizations allow the degree to grow with \(J\).

\subsection{A common return integral at the centers}

Set \(r=2s-1\) and
\[
 g(U)=\prod_{j=1}^r(U-j)=\sum_{j=0}^r b_jU^j,\qquad b_0=-r!.
\]
Every \(b_j\) is nonzero, by the signs of the elementary symmetric
functions of \(1,\ldots,r\). Begin in the real \(u,v\) plane with
\[
 \dot u=e^v-e^{2v},\qquad \dot v=g(e^u).
\]
This is \((-H_v,H_u)\) for
\(H(u,v)=G(u)+(e^v-1)^2/2\), \(G'(u)=g(e^u)\).
Each odd \(i\in\{1,3,\ldots,r\}\) gives a strict nondegenerate minimum
at \((u_i,0)=(\log i,0)\), since \(\kappa_i=G''(u_i)=i g'(i)>0\).
Choose disjoint neighborhoods of these \(s\) minima. The analytic
coordinates
\[
 w=\operatorname{sgn}(u-u_i)\sqrt{2(G(u)-G(u_i))},\qquad z=e^v-1
\]
send the energy above the minimum to \((w^2+z^2)/2\).
Write \(u=u_i+h_i(w)\), \(K_i(w)=h_i'(w)\), with
\(K_i(0)=1/\sqrt{\kappa_i}\). These changes are orientation preserving,
and the small circles are traversed counterclockwise.

Let \(t>0\) set the spatial scale. Choose real coefficients
\(a_0(t),\ldots,a_{d-1}(t)\) as specified below. For a temporarily
independent parameter \(\varepsilon\), take the analytic field
\begin{align}
 \dot u&=e^{\varepsilon a_0(t)u+v}-e^{2v},\notag\\
 \dot v&= b_0 e^{\varepsilon v/b_0}+
          \sum_{j=1}^r b_j e^{ju}
          +\varepsilon\sum_{\ell=1}^{d-1}
            \frac{a_\ell(t)}{2\ell+1}e^{(2\ell+1)v}.
 \label{cl:eq:analytic}
\end{align}
The derivative of the field with respect to \(\varepsilon\) at zero
is \((p,q)=(a_0 u e^v,\ v+\sum_{\ell=1}^{d-1}a_\ell e^{(2\ell+1)v}/(2\ell+1))\).
Its divergence is \(1+\sum_{\ell=0}^{d-1}a_\ell e^{(2\ell+1)v}\).
The two terms \(1,a_0e^v\) arise by varying exponents already present.

Let \(M_i(t,E)\) be the first variation of the energy displacement
on the oval \(w^2+z^2=t^2 E\), \(E>0\). Indeed, integrating
\(\dot H\) for one unperturbed turn gives
\(\oint p\,dv-q\,du=\iint(p_u+q_v)\,du\,dv\).
The change in return time contributes zero at first order since \(H\)
is constant along the unperturbed flow. In the circular disk the
integrand becomes \(K_i(w)\) times
\[
 \frac{1}{1+z}+\sum_{\ell=0}^{d-1}a_\ell(t)(1+z)^{2\ell}.
\]
Each vertical slice is symmetric about \(z=0\). Thus its even part
\[
 F_t(z)=\frac{1}{1-z^2}+
 \sum_{\ell=0}^{d-1}a_\ell(t)\sum_{k=0}^\ell \binom{2\ell}{2k}z^{2k}
\]
gives exactly the same integral.

Write
\[
 \mu_k=\frac{\binom{2k}{k}}{4^k(k+1)},\qquad
 \sum_{k=0}^d\mu_k c_k E^k=\mu_d\prod_{j=1}^d(E-j).
\]
All \(c_k\) are nonzero rational numbers and \(c_d=1\).
Define the \(a_\ell(t)\) by descending triangular substitution,
\begin{equation}
 1+\sum_{\ell=k}^{d-1}\binom{2\ell}{2k}a_\ell(t)
      =c_k\,t^{2(d-k)},\qquad k=d-1,\ldots,0.       \label{cl:eq:triangular}
\end{equation}
The diagonal is one. Each \(a_k\) is a polynomial in \(t\) of degree
\(2(d-k)\) with leading coefficient \(c_k\), hence is nonzero for all
sufficiently small positive \(t\). We obtain the exact identity
\begin{equation}
 F_t(z)=\sum_{k=0}^{d-1}c_k t^{2(d-k)} z^{2k}
             +\frac{z^{2d}}{1-z^2}.                 \label{cl:eq:tail}
\end{equation}

In the unit disk \(\disk=\{(\xi,\eta):\xi^2+\eta^2<1\}\), substitute
\((w,z)=t\sqrt E(\xi,\eta)\). The resulting integral is
\begin{align}
 \frac{M_i(t,E)}{t^{2d+2}}
 = E\int_{\disk} K_i(t\sqrt E\,\xi)
 \left[\sum_{k=0}^{d-1}c_k E^k\eta^{2k}
 +\frac{E^d\eta^{2d}}{1-t^2E\eta^2}\right]d\xi\,d\eta .
 \label{cl:eq:unit}
\end{align}
Polar integration gives
\(\int_{\disk}\eta^{2k}\,d\xi\,d\eta=\pi\mu_k\).
Equation \eqref{cl:eq:unit} therefore converges as \(t\to0\) to
\begin{equation}
 L_i(E)=\pi K_i(0)\mu_d\,E\prod_{j=1}^d(E-j).       \label{cl:eq:limit}
\end{equation}
On compact subintervals of \(E>0\) this convergence holds in \(C^1\)
because the bracket, \(K_i\) and their \(E\)-derivatives converge uniformly
on the fixed integration disk. The only center-dependent factor in
the limit is positive.

\subsection{From the integral to polynomial limit cycles}

We use the single choice \(\varepsilon=t^{2d+3}\). To control the return
map as the ovals shrink, first set \(w=tW,\ z=tZ,\ \varepsilon=t\sigma\)
and regard \(t,\sigma\) as independent. The scaled vector field is
analytic at \((t,\sigma)=(0,0)\) on compact annuli in \(W,Z\).
Indeed, the unscaled field in the \(w,z\) chart vanishes at
\((w,z,\varepsilon)=(0,0,0)\) for every \(t\). After the substitution,
its analytic numerator is divisible by \(t\). At \(\sigma=0\) the
scaled field is exactly
\[
 (\dot W,\dot Z)=\frac{1+tZ}{K_i(tW)}(-Z,W).
\]
On a compact annulus of positive radii containing \(E=W^2+Z^2=1,\ldots,d\)
in its interior, the angular speed stays positive for small \(t,\sigma\).
On a slightly smaller annulus initial points complete one turn, by
continuous dependence, and their return to \(Z=0,\ W>0\) is analytic in
\(E,t,\sigma\). A separate such chart is used at each of the finitely
many centers.

Write \(\Delta_i(E,t,\sigma)\) for returned minus initial \(E\).
It vanishes at \(\sigma=0\). Since \(H-H(u_i,0)=t^2E/2\),
the first variation computed above gives
\(\partial_\sigma\Delta_i(E,t,0)=2M_i(t,E)/t\) for \(t>0\).
Analyticity of the scaled return gives a Taylor remainder
\(O_{C^1(E)}(\sigma^2)\) uniform for \(t\) near zero. Setting
\(\sigma=t^{2d+2}\) yields
\begin{equation}
 \frac{\Delta_i(E,t,t^{2d+2})}{2t^{4d+3}}
     =\frac{M_i(t,E)}{t^{2d+2}}+O_{C^1(E)}(t)
     \longrightarrow L_i(E).                         \label{cl:eq:return}
\end{equation}
Choose pairwise disjoint small closed intervals about \(1,\ldots,d\)
on each of which \(L_i'\) has constant nonzero sign and the endpoint
values have opposite signs. For all sufficiently small \(t>0\) the
actual displacement has a unique simple zero in every interval, at all
\(s\) centers simultaneously. The scalar variational equation with angle
as independent variable gives a positive return derivative. Simplicity
of the zero makes this derivative different from one. These periodic
orbits are hyperbolic and isolated.
Each makes one turn and is a polar graph, so the \(d\) orbits at a given
center are distinct and nested.

To write the polynomial, take \(t=1/J\), \(\varepsilon=t^{2d+3}\),
\(\alpha=\varepsilon a_0(t)\), \(\beta=\varepsilon/b_0\), and
\[
 \Mclr{}=2J^{4d+3}\operatorname{lcm}
  \bigl(|b_0|,\operatorname{den}(c_0),\ldots,\operatorname{den}(c_{d-1})\bigr),
\]
where denominators are in lowest terms. By \eqref{cl:eq:triangular},
\(\Mclr{},\Mclr{}\alpha,\Mclr{}\beta\) are even integers. For sufficiently large \(J\),
\(|\alpha|,|\beta|<1\). In each open quadrant set \(u=\Mclr{}\log|x|\),
\(v=\Mclr{}\log|y|\). The inverse coordinate change applied to
\eqref{cl:eq:analytic}, followed by multiplication of the field by
\(\Mclr{}x^{\Mclr} y^{\Mclr}>0\) there, gives
\begin{align}
 P(x,y)={}& x^{1+\Mclr{}(1+\alpha)}y^{2\Mclr{}}-x^{1+\Mclr{}}y^{3\Mclr{}},\notag\\
 Q(x,y)={}& b_0 x^{\Mclr} y^{1+\Mclr{}(1+\beta)}
  +\sum_{j=1}^{r} b_j x^{(j+1)\Mclr{}}y^{1+\Mclr{}}
  +\varepsilon\sum_{\ell=1}^{d-1}\frac{a_\ell(t)}{2\ell+1}
                x^{\Mclr} y^{1+(2\ell+2)\Mclr{}}.                \label{cl:eq:field}
\end{align}
These polynomials define a field on the entire real plane. The exponents
in \(P\) differ in \(y\). In \(Q\), the middle sum has distinct \(x\)
exponents all greater than \(\Mclr{}\). The other terms have \(x\) exponent \(\Mclr{}\)
and pairwise different \(y\) exponents since \(-1<\beta<0\).
All coefficients are nonzero for large \(J\). Thus
\(N(X)=2+(1+r+d-1)=2s+d+1\). Each quadrant is diffeomorphic to the
whole \(u,v\) plane, and positive changes of speed preserve the section
return map. Taking the largest of the finitely many thresholds for these support conditions and for the simple return zeros proves Theorem~\ref{cl:thm:main}.

For \(s=2,d=3\), \eqref{cl:eq:triangular} gives
\((c_0,c_1,c_2)=(-15/32,55/16,-15/4)\). Thus \(L=192J^{15}\) and the cleared degree is \(1+7L\). At the algebraic sample \(J=10^{20}\), direct expansion gives eight nonzero terms and degree \(1+7\cdot192\cdot10^{300}\). The 24-cycle conclusion applies for \(J\ge J_0(2,3)\); the convergence in \eqref{cl:eq:return} leaves this threshold unquantified.

For comparison with primitive fields, set \(M=\Mclr\). The two common coordinate exponents are
\[
 a=\min\{M,1+M(1+\alpha)\},\qquad b=1+M(1+\beta).
\]
On the algebraic torus the first component vanishes precisely when \(y^M=x^{M\alpha}\). Write \(\alpha=h/k\) in lowest terms with \(k>0\). Each component of this curve has a parametrization \(x=t^k,\ y=\zeta t^h\), for an appropriate root of unity \(\zeta\). After division of \(Q\) by the unit \(x^M y^{1+M}\), its middle terms have distinct orders \(jMk\), \(1\le j\le r\). All the other orders divided by \(Mk\) are \(\alpha\beta\) or \((2\ell+1)\alpha\). For fixed \(s,d\) these have absolute value less than one at the small scales of the proof. Thus the term with \(j=r\) cannot cancel on any component. The full polynomial gcd is \(x^a y^b\). Its removal gives the same \(2s+d+1\) terms and degree \(B-a-b\). Division by the largest even monomial instead preserves the oriented returns on all open quadrants. The full quotient reverses precisely those with \(x^ay^b<0\), preserving the count of hyperbolic cycles.

